% Part: turing-machines
% Chapter: undecidability
% Section: unsolvability-decision-problem

\documentclass[../../../include/open-logic-section]{subfiles}

\begin{document}

\olfileid{tur}{und}{uns}
\olsection{निर्णय समस्या सोडवता येत नाही}

\begin{thm}
\ollabel{thm:decision-prob}
निर्णय समस्या सोडवता येत नाही:
फितीवर प्रथम-क्रम तर्कशास्त्रातील
!!a{sentence} वाक्य~$!B$ आदान म्हणून
असताना सुरू केले की $D$ अखेरीस थांबून,
$!B$ वैध असेल तेव्हाच~$1$ आणि
अन्यथा~$0$ फलित देणारे
कोणतेही ट्यूरिंग यंत्र~$D$ नाही.
\end{thm}

\begin{proof}
निर्णय समस्या सोडवता येते,
म्हणजे असे ट्यूरिंग यंत्र~$D$
आहे, असे समजा. मग पुढीलप्रमाणे
थांबण्याची समस्या सोडवता येईल.
ट्यूरिंग यंत्र~$M_e$ चा संकेतांक~$e$
आणि आदान~$w$ दिले असता
संबंधित !!{sentence} वाक्य
$!T(M_e, w) \lif !E(M_e, w)$
संगणित करून फितीवरील सर्वांत
डाव्या घराकडे पाहत थांबणारे
ट्यूरिंग यंत्र~$E$ बांधू.
मग आदान~$e$ आणि~$w$
दिले असता $E \concat D$
प्रथम $!T(M_e, w) \lif
!E(M_e, w)$ संगणित करेल,
आणि नंतर त्या आदानावर
निर्णय-समस्येचे यंत्र~$D$
चालवेल. $!T(M_e, w) \lif !E(M_e, w)$
वैध असेल तेव्हाच $D$
$1$ देऊन थांबेल, आणि
अन्यथा~$0$ देईल.
\olref[ver]{lem:halt-if-valid}
आणि \olref[ver]{lem:valid-if-halt}
यांनुसार $!T(M_e, w) \lif !E(M_e, w)$
वैध असते तेव्हाच~$M_e$
आदान~$w$ वर थांबते.
म्हणून आदान~$e$ आणि~$w$
दिले असता $E\concat D$
हे~$M_e$ आदान~$w$
वर थांबत असेल तेव्हाच
$1$ देऊन थांबेल, आणि
अन्यथा~$0$ देऊन थांबेल.
दुसऱ्या शब्दांत, $E \concat D$
थांबण्याची समस्या सोडवेल.
पण \olref[hal]{thm:halting-problem}
नुसार असे कोणतेही
ट्यूरिंग यंत्र असू शकत नाही.
\end{proof}

\begin{cor}\ollabel{cor:undecidable-sat}%
प्रथम-क्रम तर्कशास्त्रातील
कोणतेही !!{sentence} वाक्य
पूर्ततायोग्य आहे की नाही,
हे ठरवण्याची समस्या अनिर्णेय आहे.
\end{cor}

\begin{proof}
  पूर्ततायोग्यता ट्यूरिंग यंत्र~$S$
  ठरवू शकते असे समजा. मग
  निर्णय समस्या पुढीलप्रमाणे
  सोडवता येईल: आदान म्हणून
  !!a{sentence} वाक्य~$!B$ दिले
  असता~$!B$ एक घर उजवीकडे
  सरकवा. घर~$1$ कडे परत
  या आणि तेथे~$\lnot$
  हे चिन्ह लिहा.
%READERNOTE{OLTUR-016 — गोठवलेल्या इंग्रजी पुराव्यात आदान-वाक्य प्रथम B असे लिहिले आहे, पण त्याच परिच्छेदात पुढे !B सरकवले जाते आणि संपूर्ण प्रमेयात !B हेच वाक्य आहे. येथे प्रारंभीही !B लिहिले आहे; नकार आणि पूर्ततायोग्यतेची युक्ती अपरिवर्तित आहे.}

  आता ट्यूरिंग यंत्र~$S$
  चालवा. फितीवर ते अखेरीस
  $1$ ($\lnot !B$ पूर्ततायोग्य
  असेल तर) किंवा~$0$
  ($\lnot !B$ अपूर्ततायोग्य
  असेल तर) फलित देऊन
  थांबेल. घर~$1$ वर
  $\TMstroke$ असेल, तर
  तो पुसा; घर~$1$
  रिकामे असेल, तर तेथे
  $\TMstroke$ लिहा.
  त्यानंतर थांबा.

  हे ट्यूरिंग यंत्र नेहमी
  थांबते. $\lnot
  !B$ अपूर्ततायोग्य असेल
  तेव्हाच ते~$1$ देते;
  अन्यथा~$0$ देते.
  कारण $!B$ वैध असते
  तेव्हाच~$\lnot !B$
  अपूर्ततायोग्य असते,
  यंत्र~$!B$ वैध असेल
  तेव्हाच~$1$ आणि
  अन्यथा~$0$ देईल.
  म्हणजे ते निर्णय
  समस्या सोडवेल.
\end{proof}

\begin{explain}
म्हणून “$!B$ हे प्रथम-क्रम
तर्कशास्त्रातील वैध !!{sentence}
वाक्य आहे का?” या प्रश्नाला
नेहमी बरोबर “हो” किंवा
“नाही” उत्तर देणारे
ट्यूरिंग यंत्र नाही. मात्र
“हो” हे उत्तर बरोबर असेल
तेव्हा ते देणारे ट्यूरिंग
यंत्र \emph{आहे}; उत्तर
“नाही” असेल, तर ते
थांबतच नाही. हे प्रथम-क्रम
तर्कशास्त्राच्या निर्दोषता आणि
संपूर्णता प्रमेयांवरून, तसेच
!!{derivation}s प्रभावी
रीतीने प्रगणित करता येतात
यावरून मिळते.
\end{explain}

\begin{thm}
  \ollabel{thm:valid-ce}%
  प्रथम-क्रम !!{sentence}s
  वाक्यांची वैधता अर्धनिर्णेय आहे:
  फितीवर प्रथम-क्रम
  तर्कशास्त्रातील !!a{sentence}
  वाक्य~$!B$ आदान म्हणून
  असताना सुरू केले की
  $!B$ वैध असेल तेव्हाच
  अखेरीस थांबून~$1$
  देणारे, आणि अन्यथा
  न थांबणारे ट्यूरिंग
  यंत्र~$E$ आहे. हे $E$
  यंत्र वैधतेला होकार देते.
\end{thm}

\begin{proof}
  प्रथम-क्रम तर्कशास्त्रातील
  सर्व शक्य !!{derivation}s
  एका प्रभावी अल्गोरिदमने
  एकामागून एक तयार करता
  येतात. यंत्र~$E$ तसे
  करते; $\Proves
  !B$ हे दाखवणारी
  !!a{derivation} निष्पत्ती
  सापडली की ते~$1$
  देऊन थांबते. निर्दोषता
  प्रमेयामुळे $E$~हे $1$
  देऊन थांबले असेल,
  तर~$\Entails !B$.
  संपूर्णता प्रमेयामुळे
  $\Entails !B$ असेल,
  तर~$\Proves !B$
  हे दाखवणारी !!a{derivation}
  निष्पत्ती असते. $E$
  सर्व शक्य !!{derivation}s
  क्रमाने तयार करत असल्याने
  $\Proves !B$ दाखवणारी
  अशी निष्पत्ती त्याला
  अखेरीस सापडेल; म्हणून
  ते~$1$ देऊन थांबेल.
\end{proof}

\end{document}
